% DEFINED:   ch:foundations
% CONTAINS:  section/equation/figure labels live in the \input'd files below
%            (sec:nn-as-function ... sec:autograd-note, eq:residual,
%            eq:linear-grads, eq:py-binary, eq:py-softmax, eq:logodds,
%            fig:xor-boundary)
% CROSS-CHAPTER FORWARD (prose-only, no \Cref): ch.2 output heads,
%            Part II architecture, ch.8 reinforcement learning

\chapter[Foundations]{Foundations: Function Approximation and the Multilayer Perceptron}%
\label{ch:foundations}

Throughout this chapter (and the chapters that follow) we are doing
\textbf{supervised learning}: we have labeled $(x, y)$ pairs and learn $f$
such that $f(x) \approx y$. The labels can come from humans (classical
supervised) or from the data's own structure (self-supervised, as in
language modeling). Both are mathematically the same. What differs across
sections is the \emph{interpretation} of the output: the choice of
output activation (link function) and matching loss.

The network produces raw values. The output's meaning is the loss's job.

% DEFINED:   sec:nn-as-function, sec:linear-limits, eq:residual
% RESOLVED:  (none; first sections of the chapter)
% FORWARD:   sec:mlp-backprop (Task 4), sec:logistic (Task 5), sec:softmax (Task 5)

\section{What is a neural network?}\label{sec:nn-as-function}

A neural network is a parameterized function $f_\theta: \R^n \to \R^m$.
It takes a real-valued input vector and produces a real-valued output
vector. That is the whole structural commitment. What the output
\emph{means} is a separate decision, made by the loss.

The network produces raw values; the loss interprets those values as a
prediction by composing them with an \textbf{output activation} (also
called a \textbf{link function} in the language of generalized linear
models). The choice of link function and matching loss together specify
what kind of thing the network is predicting.

The simplest case is \textbf{regression}: the raw output \emph{is} the
prediction. The link function is the identity, and the loss is
\textbf{mean squared error}:
\begin{equation}\label{eq:residual}
  L = \frac{1}{2} \sum_i (\hat{y}_i - y_i)^2,
  \qquad
  \frac{\partial L}{\partial \hat{y}_i} = \hat{y}_i - y_i.
\end{equation}
The gradient is just the residual. This pattern, $\hat{y} - y$, recurs
throughout the chapter. It is the gradient of every canonical loss with
respect to its raw output, for reasons that turn out to be a theorem
about exponential families and canonical links, not an accident
(Chapter~2 proves it).

The \texttt{MSELoss} class encodes this directly:
\begin{lstlisting}
class MSELoss(Loss):
    def value(self, logits, y):
        return 0.5 * sum((zi - yi) ** 2 for zi, yi in zip(logits, y))

    def grad(self, logits, y):
        return [zi - yi for zi, yi in zip(logits, y)]

    def probs(self, logits):
        # Identity link: the prediction IS the raw output.
        return list(logits)
\end{lstlisting}
Note the $\frac{1}{2}$ convention in \texttt{value}: it makes the gradient
exactly $\hat{y} - y$ with no leading constant, matching the pattern of
the other losses.

A concrete demonstration: fit $y = \sin(x) + \varepsilon$ on
$x \in [-\pi, \pi]$ with a small MLP, identity output, and MSE loss
(\texttt{examples/regression.py}). After 2000 epochs of mini-batch SGD
the network learns the curve. This architecture is the same MLP we build
up to in \Cref{sec:mlp-backprop}; only the output head and loss differ.

With regression in hand as the unspecialized base case, the rest of the
chapter is a tour through two questions. First, what
\emph{expressive limits} does the architecture itself impose? That is
\Cref{sec:linear-limits,sec:mlp-backprop}: a single linear unit cannot
represent all continuous functions; the multilayer perceptron can.
Second, what \emph{output-side specializations} let the same
architecture solve concrete prediction problems beyond regression? That
is \Cref{sec:logistic,sec:softmax}: classification by sigmoid for
binary outcomes, softmax for categorical.

\section{The single linear unit and its limits}\label{sec:linear-limits}

The smallest non-trivial network is one \texttt{Linear} layer mapping
inputs to outputs:
\[
  \vz = W \vx + \mathbf{b}.
\]
This is a \textbf{single-layer perceptron} (SLP): an affine map with no
non-linearity. Trained with MSE it learns the least-squares fit. Trained
with a sigmoid head and BCE loss (\Cref{sec:logistic}), it is logistic
regression. The function class either way is \emph{linear} in the input.

The SLP can fit any linearly separable target. It cannot fit anything
else.

The canonical demonstration is XOR on $\{0, 1\}^2$:
\[
  \mathrm{XOR}(0,0) = 0, \quad
  \mathrm{XOR}(0,1) = 1, \quad
  \mathrm{XOR}(1,0) = 1, \quad
  \mathrm{XOR}(1,1) = 0.
\]
Training a single \texttt{Linear} layer with \texttt{SigmoidBCE} on these
four points for 500 epochs (batch size 4, seed~42):
\begin{lstlisting}
slp = Network([Linear(2, 1)], loss=SigmoidBCE())
slp.fit(xor_X, xor_Y, epochs=500, lr=0.5, batch_size=4)
\end{lstlisting}
Every input converges to $p = 0.5000$. The model literally cannot do
better. To see why, write down what fitting all four points would
require. With $\hat{y} = \sigma(w_1 x_1 + w_2 x_2 + b)$ and targets
$\{0, 1, 1, 0\}$, the BCE loss is minimized when each $\hat{y}$ matches
its target. Since sigmoid is monotone, the four constraints reduce to
four inequalities:
\[
  b < 0, \quad
  w_1 + b > 0, \quad
  w_2 + b > 0, \quad
  w_1 + w_2 + b < 0.
\]
Adding the second and third gives $w_1 + w_2 + 2b > 0$, so
$w_1 + w_2 + b > -b > 0$. But the fourth requires
$w_1 + w_2 + b < 0$. Contradiction. No choice of $(w_1, w_2, b)$
fits all four points simultaneously.

Geometrically: a linear classifier draws a single hyperplane in the
input space (a line in $\R^2$). XOR's positive class
$\{(0,1),(1,0)\}$ and negative class $\{(0,0),(1,1)\}$ lie at
diagonal corners of the unit square. No line separates them.

This was \textcite{minsky1969perceptrons}'s argument in 1969, which
stalled the perceptron research program for a decade and a half. The
resolution is one of the cleanest examples of how architecture and
expressivity interact. To escape linearity, add a non-linearity: the
subject of \Cref{sec:mlp-backprop}.

% DEFINED:   sec:mlp-backprop, eq:linear-grads
% RESOLVED:  sec:linear-limits, sec:nn-as-function
% FORWARD:   (none; intra-chapter refs resolved after Task 4)

\section{The multilayer perceptron, and what depth and non-linearity buy}\label{sec:mlp-backprop}

% TODO Task 4

% DEFINED:   sec:logistic, sec:softmax, eq:py-binary, eq:py-softmax
% RESOLVED:  sec:linear-limits, sec:mlp-backprop, sec:nn-as-function
% FORWARD:   sec:inductive-biases (Task 7)

\section{Logistic regression}\label{sec:logistic}

% TODO Task 5

\section{Multi-class regression via softmax}\label{sec:softmax}

% TODO Task 5

% DEFINED:   sec:logits-logprobs, eq:logodds
% RESOLVED:  sec:softmax
% FORWARD:   sec:stability (Task 8)

\section{Logits are log-probabilities, up to a constant}\label{sec:logits-logprobs}

% TODO Task 6

% DEFINED:   sec:inductive-biases
% RESOLVED:  sec:softmax
% FORWARD REFS (prose-only): Chapter 2 (output-head axis),
%             Part II (architecture axis), Chapter 8 (RL)

\section{Inductive biases}\label{sec:inductive-biases}

% TODO Task 7

% DEFINED:   sec:stability, sec:gradcheck, sec:autograd-note
% RESOLVED:  sec:logits-logprobs, sec:mlp-backprop
% FORWARD:   (none)

\section{Numerical stability is the gauge freedom, reused}\label{sec:stability}

$e^z$ overflows for moderately large $z$. The fix for softmax: subtract
the largest logit first,
\[
  \softmax(\vz) = \softmax\!\left(\vz - \max_i z_i\right).
\]
This is \emph{exactly} the additive-constant freedom from
\Cref{sec:logits-logprobs}. Softmax does not care about a shared shift
of the logits, so we spend that freedom to keep the exponentials in a
safe range. Stability here is not a new idea bolted on. It is the gauge
symmetry put to work.

Cross-entropy is computed as $\mathrm{logsumexp}(\vz) - z_y$ rather than
$-\log \mathrm{softmax}(\vz)_y$, which would risk $\log 0$. The sigmoid
is evaluated with a sign branch; binary cross-entropy is computed straight
from the logit, both to avoid overflow and $\log 0$. These are the
primitives the loss classes call:

\begin{lstlisting}
def sigmoid(z):
    if z >= 0.0:
        return 1.0 / (1.0 + math.exp(-z))     # safe: exp(-z), z >= 0
    ez = math.exp(z)                          # safe: exp(z),  z < 0
    return ez / (1.0 + ez)

def logsumexp(zs):
    m = max(zs)
    return m + math.log(sum(math.exp(z - m) for z in zs))

def softmax(zs):
    m = max(zs)
    exps = [math.exp(z - m) for z in zs]
    total = sum(exps)
    return [e / total for e in exps]
\end{lstlisting}

In every case the exponent argument is held at or below $0$, so $e^z$
stays in $[0, 1]$ and never overflows. The max-subtraction in
\texttt{softmax} and \texttt{logsumexp} is the \Cref{sec:logits-logprobs}
gauge freedom; the sign branch in \texttt{sigmoid} is the same trick for
the two-class case.

\section{Trust, but verify: numerical gradients}\label{sec:gradcheck}

Hand-derived gradients are easy to get subtly wrong. There is a slow but
foolproof check: the definition of a derivative itself. For any parameter
$\theta$,
\[
  \frac{\partial L}{\partial \theta}
  \approx
  \frac{L(\theta + \varepsilon) - L(\theta - \varepsilon)}{2\varepsilon}.
\]
\texttt{gradient\_check} computes this central difference for every
parameter and compares it to the analytical gradient from
\texttt{backward}. It perturbs each parameter in place, twice, and reads
off the loss:

\begin{lstlisting}
def gradient_check(net, x, y, eps=1e-5):
    net.zero_grad()
    net.forward(x)
    net.backward(y)
    worst = 0.0
    for values, grads in net.parameters():
        for k in range(len(values)):
            original = values[k]
            values[k] = original + eps
            loss_plus = net.loss_value(x, y)
            values[k] = original - eps
            loss_minus = net.loss_value(x, y)
            values[k] = original
            numerical = (loss_plus - loss_minus) / (2.0 * eps)
            analytical = grads[k]
            denom = max(abs(numerical) + abs(analytical), 1e-12)
            worst = max(worst, abs(numerical - analytical) / denom)
    return worst
\end{lstlisting}

This is too slow to train with: one forward pass per parameter. But it is
the ground truth that keeps the fast method honest. It works for
\emph{any} layer and loss because it only ever touches
\texttt{parameters()}, \texttt{forward}, and \texttt{backward}. Running
it on the four configurations from \Cref{sec:mlp-backprop} (seed~42):

\begin{itemize}
\item Logistic regression (SigmoidBCE): worst relative error $\approx 6 \times 10^{-12}$
\item Softmax regression (SoftmaxCE): $\approx 1 \times 10^{-10}$
\item Tanh MLP (SoftmaxCE): $\approx 7 \times 10^{-9}$
\item ReLU MLP (SigmoidBCE): $\approx 1 \times 10^{-9}$
\end{itemize}

All residuals are near machine precision. One caveat: \texttt{ReLU} has a
kink at $0$ and no derivative there. A finite difference that straddles
the kink will disagree with the analytical sub-gradient. With random
initial weights the probability of a pre-activation landing exactly on
the kink is effectively zero, so the check passes in practice. The
tolerance is anchored on a \texttt{Tanh} network, which is smooth
everywhere; the $10^{-9}$ figure for the ReLU network is real, not a
relaxed threshold masking a bug.

\section{Closing note: from per-layer to per-operation}\label{sec:autograd-note}

Backpropagation here is per-\emph{layer}: each layer carries a local
\texttt{backward}. But nothing forced the unit to be a layer. If every
scalar \emph{operation} (every $+$, every $\times$, every $\tanh$)
carried its own local backward, the same reverse pass would compute
gradients through an arbitrary expression, with no hand-derived layer
gradients at all. That is \textbf{automatic differentiation}, and it is
what every modern framework's autograd engine does
(\cite{baydin2018autodiff}). It is the same idea as this library, at a
finer grain. Building one is the natural next step, and Chapter~2
returns to this template when framing its own natural extensions.

% DEFINED:   (none; unnumbered section)
% RESOLVED:  \cite keys only
% FORWARD:   (none)

\section*{Bibliographic Notes}

% TODO Task 9

